\subsection{Escalamiento}

El proceso de escalamiento garantiza que las incidencias relacionadas con el servicio de red sean atendidas de manera oportuna y por el nivel de soporte adecuado, de acuerdo con la estructura de contactos definida en este SLA.

\subsubsection*{Niveles de Escalamiento}
\begin{itemize}
    \item \textbf{Nivel 1: Mesa de Ayuda de Aulas.} Primer punto de contacto para el registro de incidentes y solicitudes. Realiza el diagnóstico inicial, aplica soluciones estándar y da seguimiento básico. Si el incidente no puede ser resuelto o se detecta que afecta a múltiples usuarios o servicios críticos, se escalara al Nivel 2 con toda la información recolectada (descripción, impacto, bitácora de acciones y evidencias).
    \item \textbf{Nivel 2: Soporte Especializado de Redes.} Atiende incidentes escalados desde el Nivel 1 que involucren conmutadores, enlaces troncales, VLANs, Wi-Fi o firewall perimetral. Si el incidente presenta impacto crítico (por ejemplo, indisponibilidad total del edificio o riesgo de seguridad), o requiere cambios mayores en la infraestructura, el caso se escalara al Nivel 3.
    \item \textbf{Nivel 3: Coordinación de Infraestructura y Proveedores.} Responsable de la gestión de incidentes mayores, coordinación de cambios significativos, relación con el proveedor de Internet y fabricantes, así como de la toma de decisiones para la restauración del servicio cuando se requieran acciones extraordinarias.
\end{itemize}

En cada cambio de nivel, se actualiza el estado del incidente en la bitácora correspondiente y se notifica a la parte contratante sobre el progreso y las acciones en curso.

\subsubsection*{Clasificación de Incidencias y Tiempos de Atención}
Las incidencias se clasifican según su impacto y urgencia, definiendo tiempos objetivo de respuesta (inicio de la atención) y de resolución (restauración del servicio o aplicación de una solución temporal aceptable). La Tabla~\ref{tab:escalamiento-incidencias} resume estos parámetros.

\begin{table}[H]
    \centering
    \caption{Clasificación de incidencias y tiempos de atención}\label{tab:escalamiento-incidencias}
    \begin{tabular}{p{4cm}p{5cm}cc}
        	\toprule
        	\textbf{Clasificación} & \textbf{Descripción} & \textbf{Tiempo de respuesta} & \textbf{Tiempo de resolución} \\
        \midrule
        Crítica & Caída total del servicio de red en el edificio o incidente de seguridad grave. & Hasta 2 hora & Hasta 4 horas \\
        Alta & Afectación parcial significativa (varias aulas o servicios clave). & Hasta 2 horas & Hasta 8 horas \\
        Media & Problemas que afectan a un conjunto reducido de usuarios o servicios no críticos. & Hasta 2 horas & Hasta 24 horas \\
        Baja & Solicitudes de servicio, dudas o mejoras sin impacto directo en la disponibilidad. & Hasta 2 día hábil & Hasta 5 días hábiles \\
        \bottomrule
    \end{tabular}
\end{table}
